\documentclass[11pt,letterpaper]{article}
\usepackage[utf8]{inputenc}
\usepackage[T1]{fontenc}
\usepackage[spanish,es-tabla]{babel}
\usepackage{mathptmx}
\usepackage{graphicx}
\usepackage{array}
\usepackage{longtable}
\usepackage{geometry}
\usepackage{fancyhdr}
\usepackage{cite}
\usepackage[hidelinks,linktoc=all]{hyperref}
\geometry{top=2.3cm,bottom=2.3cm,left=2.5cm,right=2.2cm,headheight=14pt,headsep=16pt,footskip=30pt}
\fancyhf{}
\fancyhead[L]{\fontsize{9}{11}\selectfont\textbf{Pontificia Universidad Javeriana}}
\fancyhead[R]{\fontsize{9}{11}\selectfont Propuesta para Trabajo de Grado}
\fancyfoot[R]{\fontsize{9}{11}\selectfont Página \thepage}
\renewcommand{\headrulewidth}{0.4pt}
\renewcommand{\footrulewidth}{0.4pt}
\pagestyle{fancy}
\fancypagestyle{plain}{\fancyhf{}\fancyhead[L]{\fontsize{9}{11}\selectfont\textbf{Pontificia Universidad Javeriana}}\fancyhead[R]{\fontsize{9}{11}\selectfont Propuesta para Trabajo de Grado}\fancyfoot[R]{\fontsize{9}{11}\selectfont Página \thepage}\renewcommand{\headrulewidth}{0.4pt}\renewcommand{\footrulewidth}{0.4pt}}
\setlength{\parindent}{0pt}
\setlength{\parskip}{7pt}
\setlength{\emergencystretch}{2em}
\clubpenalty=10000
\widowpenalty=10000
\brokenpenalty=10000
\renewcommand{\baselinestretch}{1.08}
\setcounter{secnumdepth}{3}
\setcounter{tocdepth}{3}
\hypersetup{pdftitle={Propuesta de proyecto de grado: apoyo a la decisión para incendios},pdfauthor={Andrés Pinzón, Daniel Ramírez, Guillermo Aponte, Miguel Sarmiento}}
\newcommand{\datosEstudiante}[1]{\begin{minipage}{\linewidth}\textbf{#1}\enspace\hrulefill\par\vspace{1mm}{\fontsize{8.5}{10}\selectfont\begin{tabular}[t]{@{}p{0.17\linewidth}p{0.16\linewidth}p{0.23\linewidth}p{0.34\linewidth}@{}}\textbf{Documento}&\textbf{Celular}&\textbf{Teléfono fijo}&\textbf{Correo Javeriano}\\Por registrar&Por registrar&Por registrar&Por registrar\end{tabular}\par}\end{minipage}}
\begin{document}
\clearpage
\newgeometry{top=2.3cm,bottom=2cm,left=2cm,right=1.5cm}
\setlength{\headwidth}{\textwidth}
\thispagestyle{fancy}
\noindent\begin{minipage}[b]{0.52\linewidth}\includegraphics[height=2.2cm]{logo-javeriana.png}\end{minipage}\hfill\begin{minipage}[b]{0.43\linewidth}\raggedleft\fontsize{14}{17}\selectfont Ingeniería de Sistemas\end{minipage}
\vspace{0.8cm}
\begin{center}\fontsize{16}{19}\selectfont\bfseries PROPUESTA PARA PROYECTO DE GRADO\end{center}
\vspace{-0.4cm}
\begingroup\setlength{\tabcolsep}{3pt}\setlength{\arrayrulewidth}{0.8pt}\fontsize{10}{12}\selectfont
\noindent\begin{tabular}{|l|}\hline
\begin{minipage}{\dimexpr\linewidth-2\tabcolsep-2\arrayrulewidth\relax}\vspace{2mm}\raggedright
\textbf{TÍTULO}\par\vspace{1mm}
Aplicación de apoyo a la decisión para incendios en las zonas urbanas y forestales de los Cerros Orientales de Bogotá\par\vspace{3mm}
\end{minipage}\\\hline
\begin{minipage}{\dimexpr\linewidth-2\tabcolsep-2\arrayrulewidth\relax}\vspace{2mm}\raggedright
\textbf{OBJETIVO GENERAL}\par\vspace{1mm}
Desarrollar una aplicación de escritorio para estimar y visualizar el comportamiento probable de incendios forestales en un sector de los Cerros Orientales de Bogotá, a partir de escenarios WFDS precalculados, con el fin de apoyar al comandante o coordinador de la operación.\par\vspace{3mm}
\end{minipage}\\\hline
\begin{minipage}{\dimexpr\linewidth-2\tabcolsep-2\arrayrulewidth\relax}\vspace{2mm}\raggedright
\textbf{ESTUDIANTE(S) [Pendientes: 1. completar datos de identificación y contacto]}\par\vspace{2mm}
\datosEstudiante{Andrés Pinzón}\par\vspace{3mm}
\datosEstudiante{Daniel Ramírez}\par\vspace{3mm}
\datosEstudiante{Guillermo Aponte}\par\vspace{3mm}
\datosEstudiante{Miguel Sarmiento}\par\vspace{4mm}
\end{minipage}\\\hline
\begin{minipage}{\dimexpr\linewidth-2\tabcolsep-2\arrayrulewidth\relax}\vspace{2mm}\raggedright
\textbf{DIRECTOR [Pendientes: 1. completar datos de identificación y contacto]}\par\vspace{2mm}
\textbf{Jaime Andrés Pavlich Mariscal}\enspace\hrulefill\par\vspace{1mm}
{\fontsize{8}{9.5}\selectfont\begin{tabular}[t]{@{}p{0.13\linewidth}p{0.13\linewidth}p{0.15\linewidth}p{0.22\linewidth}p{0.27\linewidth}@{}}
\textbf{Documento} & \textbf{Celular} & \textbf{Teléfono fijo} & \textbf{Correo Javeriano} & \textbf{Empresa donde trabaja y cargo}\\
Por registrar & Por registrar & Por registrar & Por registrar & Pontificia Universidad Javeriana. Cargo por registrar.
\end{tabular}\par}\vspace{4mm}
\end{minipage}\\\hline
\end{tabular}\endgroup
\vfill\begin{flushright}\small 26 de octubre de 2026\end{flushright}
\restoregeometry
\setlength{\headwidth}{\textwidth}
\clearpage

\clearpage
\thispagestyle{fancy}
\noindent\begin{minipage}[b]{0.52\linewidth}\includegraphics[height=2.2cm]{logo-javeriana.png}\end{minipage}\hfill\begin{minipage}[b]{0.43\linewidth}\raggedleft\fontsize{14}{17}\selectfont Ingeniería de Sistemas\end{minipage}
\vspace{5mm}
\noindent\textbf{Historial de Cambios}\par\vspace{2mm}
\begingroup\small\setlength{\tabcolsep}{5pt}\renewcommand{\arraystretch}{1.2}
\noindent\begin{tabular}{|p{0.12\linewidth}|p{0.17\linewidth}|p{\dimexpr0.71\linewidth-6\tabcolsep-4\arrayrulewidth\relax}|}\hline
Fecha & Autor & Descripción\\[2mm]
25-Oct-2026 & Daniel Felipe Ramirez Vargas & Primera versión. Incorpora el contexto vigente del proyecto, las especificaciones de la plantilla, los aportes pertinentes que hemos realizado en trabajos pasados como el BOI, VPP y PAFP, y las indicaciones de la profesora y del director. Queda registrado en cada sección que aspectos siguen pendientes por resolver y se utiliza citación IEEE.\\\hline
27/10/2026 & Daniel Felipe Ramirez Vargas & Incorporación de glosarios, actualización de ficha de estudiantes y unificación de terminología FDS/WFDS.\\\hline
\end{tabular}\endgroup
\clearpage

\tableofcontents
\clearpage
\section{Visión global}\label{sec:vision}
\subsection{Antecedentes, problema y solución propuesta}
\subsubsection{Descripción de la problemática u oportunidad}\label{sec:problematica}
Los incendios forestales de los Cerros Orientales de Bogotá plantean un problema de interpretación del territorio durante la respuesta, debido a que el avance del fuego depende de condiciones de viento, relieve y vegetación que pueden variar mientras se organiza la atención. En las zonas próximas a la ciudad, esa evolución puede comprometer viviendas, vías y otros elementos expuestos, de modo que conocer la posición observada del incendio resulta necesario, pero no basta para explorar hacia dónde podría propagarse. Las entrevistas realizadas por el equipo identifican necesidades relacionadas con la dirección y velocidad del avance, los accesos, las posibles barreras y la exposición de personas e infraestructura, además de dificultades para recibir y verificar información oportuna. Estas observaciones se utilizan como evidencia cualitativa de los participantes y no como una afirmación de que todas las instituciones carecen de herramientas de apoyo. En particular, el comandante o coordinador de la operación necesita interpretar información cambiante para orientar su juicio, por lo cual la oportunidad del proyecto consiste en acercar resultados de simulación a una consulta comprensible dentro de ese contexto de decisión~\cite{entrevistas2026}.

La dificultad también tiene una dimensión técnica, ya que representar físicamente la interacción entre fuego, combustible y atmósfera demanda datos y capacidad de cálculo que no pueden suponerse disponibles durante un incidente. Los antecedentes sobre los Cerros Orientales muestran, además, que la caracterización territorial y la calibración de combustibles son problemas propios del contexto local, por lo que disponer de un simulador no resuelve automáticamente la preparación de sus entradas~\cite{verano2013,mendez2018}. De esta manera, existe una oportunidad de Ingeniería de Sistemas en la organización de datos, resultados y mecanismos de consulta que permitan aprovechar preparación computacional previa, sin trasladar al usuario operativo la configuración detallada del modelo ni exigirle interpretar directamente sus archivos científicos. La pertinencia de esa integración fue reforzada en la retroalimentación académica del proyecto, que solicita mantener la aplicación como producto central y hacer explícitas sus condiciones de viabilidad~\cite{pafp2026}.

\subsubsection{Formulación del problema [Pendientes: 1. delimitar el sector de estudio]}\label{sec:problema}
¿Cómo ofrecer al comandante o coordinador de la operación una estimación consultable y una visualización georreferenciada de la posible propagación de un incendio forestal en un sector de los Cerros Orientales de Bogotá, aprovechando resultados calculados previamente y reduciendo la necesidad de ejecutar simulaciones físicas durante el incidente? La pregunta se concentra en transformar resultados científicos en información accesible para el usuario descrito en la sección~\ref{sec:problematica}, dentro de las condiciones representadas por una biblioteca finita. Por ello, la solución deberá permitir reconocer tanto los escenarios compatibles con la situación observada como sus límites de representación, sin confundir una estimación condicionada por datos y supuestos con una predicción cierta del comportamiento real. El ámbito territorial definitivo permanece pendiente y su selección determinará la extensión de los datos, los casos de simulación y las comprobaciones que podrán sostenerse en el proyecto.

\subsubsection{Propuesta de solución [Pendientes: 1. concretar el sector y la cobertura de escenarios, 2. definir selección y agregación, 3. fijar horizonte de consulta y servicio meteorológico]}\label{sec:solucion}
Se propone desarrollar una aplicación de escritorio sustentada en una biblioteca de escenarios precalculados con WFDS. Su preparación comprenderá la configuración del simulador, la elaboración de entradas territoriales compatibles y la ejecución anticipada de casos mediante recursos de cómputo de alto desempeño, de acuerdo con las capacidades que se comprueben. Las igniciones de referencia se distribuirán con densidad variable guiada por amenaza de ignición, de manera que las zonas de mayor propensión tengan más puntos y las de menor propensión conserven cobertura con puntos más separados. Esta distribución corresponde a una decisión del proyecto, cuyo antecedente territorial es el trabajo de Morales Salguero~\cite{morales2017}, y debe distinguirse de la discretización tridimensional utilizada para resolver cada simulación. Los resultados se almacenarán con su ubicación, condiciones de entrada, tiempos y procedencia, de modo que la aplicación pueda recuperar escenarios sin ejecutar de nuevo el cálculo físico completo durante una consulta.

En el uso previsto, el usuario proporcionará una ubicación inicial o un perímetro observado y la aplicación incorporará condiciones meteorológicas actuales obtenidas mediante una API. A partir de esos datos se seleccionarán escenarios próximos espacialmente y compatibles con las condiciones representadas en la biblioteca. Cuando exista un perímetro, la consulta considerará igniciones cercanas al borde activo y agregará sus resultados espaciales, excluyendo el interior que ya se identifica como quemado. La salida se mostrará como una animación georreferenciada para un horizonte limitado, denominado $\Delta t$, y podrá actualizarse ante nuevas observaciones o cambios relevantes del viento. La forma de obtener y verificar el perímetro, las reglas de semejanza, las ponderaciones y la estabilidad temporal que determine $\Delta t$ deberán precisarse y comprobarse durante el desarrollo. En particular, combinar resultados de igniciones independientes constituye una hipótesis de aproximación que requiere contraste, por lo que la coincidencia espacial entre escenarios no se presentará como una probabilidad calibrada de afectación.

La aplicación apoyará la interpretación del comandante o coordinador y permitirá explorar los escenarios como material de capacitación, sin asumir decisiones operativas automáticas ni incorporar un simulador táctico de colocación de recursos. El alcance tampoco incluye por ahora recepción de reportes comunitarios, gestión de recursos, integración de sensores o interoperabilidad institucional, aunque estas necesidades se conservarán como hallazgos de las entrevistas. La preparación anticipada reduce la necesidad de simulación durante el incidente, pero no acredita funcionamiento completo sin conexión, pues la consulta meteorológica y las condiciones de contingencia todavía requieren definición. El diseño anterior del póster~\cite{poster2026} sirve como antecedente documental, mientras que esta propuesta recoge las decisiones posteriores de escritorio, WFDS exclusivo y densidad espacial variable.

\subsubsection{Justificación de la solución [Pendientes: 1. sustentar la comparación con las alternativas de la sección 5.2]}\label{sec:justificacion}
La solución se justifica por la posibilidad de separar el esfuerzo de preparación científica del momento en que se necesita consultar información para decidir. En lugar de exigir al usuario configurar modelos y esperar sus ejecuciones, el producto busca convertir resultados previamente organizados en una representación interpretable de la evolución posible del incendio. El valor esperado se encuentra en esa articulación de biblioteca, consulta y comunicación espacial, mientras que sus fundamentos computacionales se desarrollan en la sección~\ref{sec:fundamentos}. Este enfoque resulta pertinente para la necesidad expresada en las entrevistas, aunque su utilidad depende de que las condiciones recuperadas representen suficientemente la situación consultada y de que la interfaz permita entender sus limitaciones. En consecuencia, la propuesta incluye comprobaciones técnicas y evaluación de comprensión, en vez de atribuir por anticipado exactitud física o efectividad operativa a una animación.

La elección también busca acotar un problema territorial y computacional amplio a una contribución realizable de software, con preparación trazable de datos y resultados. Sin embargo, este argumento de pertinencia aún no constituye una comparación completa de soluciones, dado que las tres alternativas adicionales exigidas por la plantilla no se han establecido. La justificación comparativa se completará a partir del análisis de la sección~\ref{sec:alternativas}, utilizando criterios comunes de alcance, datos, costo de preparación, consulta y presentación de resultados. Hasta entonces, no se afirma que la propuesta supere a herramientas existentes ni que las tesis empleadas como antecedentes equivalgan, por sí mismas, a alternativas integrales para el mismo problema.


\subsection{Descripción general del proyecto}\label{sec:descripcion}
El proyecto se plantea como un desarrollo aplicado de Ingeniería de Sistemas cuyo producto central será la aplicación, mientras que la preparación de simulaciones constituirá el soporte de sus consultas. Su alcance funcional y territorial se establece en la sección~\ref{sec:solucion}, y los objetivos siguientes expresan la contribución que se espera demostrar mediante los entregables de la sección~\ref{sec:entregables}.
\subsubsection{Objetivo general [Pendientes: 1. precisar sector, 2. acordar criterio global de cumplimiento y plazo]}\label{sec:objetivogeneral}
Desarrollar una aplicación de escritorio para estimar y visualizar el comportamiento probable de incendios forestales en un sector de los Cerros Orientales de Bogotá, a partir de escenarios FDS precalculados, con el fin de apoyar al comandante o coordinador de la operación. El cumplimiento se documentará mediante una versión integrada que permita demostrar el flujo descrito en la sección~\ref{sec:solucion}, junto con evidencia de pruebas y evaluación de comprensión. La delimitación territorial, las condiciones concretas de esa demostración y el compromiso temporal deberán acordarse para cerrar la formulación verificable del objetivo, sin anticipar resultados de precisión o desempeño que todavía no se han medido.

\subsubsection{Objetivos específicos [Pendientes: 1. acordar métricas y umbrales, 2. precisar tamaño de biblioteca, 3. confirmar plazos de cumplimiento]}\label{sec:objetivos}
Los cuatro objetivos del póster se conservan como base y se ajustan en su redacción al alcance vigente de software. La retroalimentación del PAFP solicita reforzar ese enfoque, pero no contiene una lista definitiva que sustituya los objetivos previos~\cite{pafp2026,poster2026}. Por esta razón, se propone la siguiente formulación de trabajo, cuya medición se concretará con los criterios de las fases y los valores pendientes indicados en este título.
\begin{enumerate}
\item \textbf{OE1.} Diseñar e implementar un mecanismo de búsqueda que recupere escenarios de la biblioteca a partir de la ubicación o el perímetro observado y de condiciones meteorológicas actuales obtenidas mediante una API, conservando la identificación de los casos seleccionados para comprobar la correspondencia entre la consulta y sus resultados.
\item \textbf{OE2.} Construir una biblioteca consultable de simulaciones WFDS precalculadas para las igniciones y condiciones acordadas del sector de estudio, mediante ejecuciones en HPC y las modalidades de paralelización que se verifiquen como compatibles, con metadatos que permitan identificar entradas, configuración y salidas de cada caso.
\item \textbf{OE3.} Configurar WFDS con niveles de complejidad de parámetros pertinentes para representar la propagación en las coberturas urbanas y forestales presentes en el sector, preparando entradas reproducibles y documentando sus restricciones, de manera que sus resultados puedan incorporarse al flujo de la aplicación.
\item \textbf{OE4.} Desarrollar e integrar la aplicación de escritorio con animaciones georreferenciadas del avance estimado del fuego y actualización de consultas, para que sus resultados y limitaciones puedan ser interpretados por el usuario previsto sin formación especializada en modelado, mediante tareas de evaluación que aporten evidencia de esa comprensión.
\end{enumerate}


\clearpage
\subsection{Entregables, estándares utilizados y justificación [Pendientes: 1. acordar criterios cuantitativos de aceptación]}\label{sec:entregables}
Los entregables se identificarán por versión y se relacionarán con los objetivos y las evidencias de cumplimiento, de modo que puedan ser revisados durante las fases de la sección~\ref{sec:proceso}. La tabla~\ref{tab:entregables} concentra los productos comprometidos y el propósito de los estándares seleccionados. Su utilización será una adaptación académica explícita al alcance del proyecto, por lo que citar una norma no implica certificación ni demuestra conformidad completa. Asimismo, los estándares de ingeniería y calidad apoyarán la organización de la evidencia, mientras que la validez de datos, simulaciones y agregación deberá sustentarse mediante comprobaciones técnicas específicas.

{\small\setlength{\tabcolsep}{5pt}\renewcommand{\arraystretch}{1.18}
\begin{longtable}{|>{\raggedright\arraybackslash}p{0.26\linewidth}|>{\raggedright\arraybackslash}p{0.28\linewidth}|>{\raggedright\arraybackslash}p{0.36\linewidth}|}
\caption{Entregables y aplicación de estándares.}\label{tab:entregables}\\\hline
\textbf{Entregable} & \textbf{Estándares asociados} & \textbf{Justificación}\\\hline\endfirsthead
\hline\textbf{Entregable} & \textbf{Estándares asociados} & \textbf{Justificación}\\\hline\endhead
E1. Especificación de requisitos y matriz de trazabilidad. & ISO/IEC/IEEE 29148:2018~\cite{iso29148}. & Relacionar necesidades, alcance, requisitos verificables y evidencia, conservando la procedencia de las necesidades y las exclusiones acordadas.\\\hline
E2. Documento de arquitectura y diseño. & OMG UML 2.5.1~\cite{uml251} e ISO/IEC/IEEE 12207:2026~\cite{iso12207}. & Describir componentes, interfaces y flujo de información, vinculando el diseño con el desarrollo e integración de la solución.\\\hline
E3. Paquete reproducible de configuración WFDS y preparación de datos. & ISO/IEC/IEEE 12207:2026 e IEEE 730:2026~\cite{ieee730}. & Identificar versiones, entradas, unidades, procedencia y evidencia de revisión de los productos técnicos.\\\hline
E4. Biblioteca de escenarios y catálogo de metadatos. & ISO/IEC/IEEE 12207:2026 e IEEE 730:2026. & Conservar la relación entre cada resultado, sus condiciones y su configuración, para permitir consulta e identificar casos incompletos o incompatibles.\\\hline
E5. Código fuente y versión ejecutable de la aplicación de escritorio. & ISO/IEC/IEEE 12207:2026 y OMG UML 2.5.1. & Mantener correspondencia entre diseño, implementación e integración de los componentes de consulta y visualización.\\\hline
E6. Especificaciones, registros e informe de pruebas y evaluación. & ISO/IEC/IEEE 29119-2, 29119-3 y 29119-4:2021~\cite{iso291192,iso291193,iso291194}, e IEEE 730:2026. & Organizar procesos, documentación y técnicas de prueba, registrando condiciones, resultados y reevaluaciones frente a criterios explícitos.\\\hline
E7. Manual de instalación y uso e informe final del proyecto. & ISO/IEC/IEEE 12207:2026 e IEEE 730:2026. & Acompañar la entrega con instrucciones reproducibles, dependencias, límites conocidos y síntesis de la evidencia obtenida.\\\hline
\end{longtable}}


\clearpage
\section{Análisis de impacto}\label{sec:impacto}
A corto plazo, el impacto esperado se concentra en disponer de una base técnica organizada y una herramienta demostrable que acerque al usuario previsto a escenarios de comportamiento del fuego. La contribución inmediata será facilitar la lectura conjunta de la situación consultada, su evolución estimada y las condiciones bajo las cuales se obtuvo el resultado, en lugar de entregar únicamente archivos de simulación. La evidencia de este aporte deberá provenir de las comprobaciones y tareas de evaluación descritas en la sección~\ref{sec:transicion}, por lo que no se presume una mejora operativa antes de realizarlas. Para el equipo académico, el desarrollo permitirá consolidar una forma reproducible de preparar datos, identificar restricciones del modelo y conectar productos científicos con requisitos de una aplicación, lo cual constituye un resultado de formación e integración tecnológica dentro del alcance definido.

A mediano plazo, la disponibilidad de escenarios consultables podría servir como apoyo a la capacitación y a la discusión de situaciones de propagación con personas vinculadas a la atención de incendios. Su utilidad social esperada radica en favorecer una comprensión compartida de la exposición y de los límites de la estimación, especialmente cuando las decisiones afectan comunidades próximas al área estudiada. Esta contribución dependerá de que los participantes comprendan la representación y puedan reconocer situaciones que la biblioteca no cubre, de modo que el producto deberá comunicar las restricciones junto con sus resultados. No se compromete una adopción institucional, una sustitución de procedimientos de respuesta ni una disminución cuantificada de tiempos, daños o pérdidas, ya que ninguna de esas consecuencias ha sido evaluada por el proyecto.

A largo plazo, el trabajo podría aportar conocimiento reutilizable sobre las condiciones necesarias para convertir simulaciones físicas en apoyo informativo para la gestión de incendios. En particular, el registro de datos faltantes, resultados desfavorables y límites de aproximación puede orientar a quienes consulten el proyecto y evitar que una representación visual convincente se interprete como certeza científica. Este aporte se relaciona temáticamente con la protección de ecosistemas terrestres y la acción frente a condiciones climáticas consideradas en el póster~\cite{poster2026}, aunque no equivale a demostrar una contribución medible a indicadores de desarrollo sostenible. La proyección se refiere a la utilidad del conocimiento y del producto entregado, sin ampliar el compromiso territorial ni establecer servicios permanentes que excedan la propuesta.

\clearpage
\section{Proceso}\label{sec:proceso}
Se utilizará Disciplined Agile Delivery, DAD, organizado en Inicio, Construcción y Transición, conforme a la metodología adoptada por el proyecto~\cite{poster2026}. La adaptación considera que la aplicación depende tanto de la ingeniería de software como de preparación de datos y simulación, por lo que las fases incorporan resultados técnicos y funcionales revisables. Los objetivos no se asignan de forma exclusiva a una sola fase, ya que cada uno recibirá aportes sucesivos hasta su comprobación integrada. Los entregables se identifican en la sección~\ref{sec:entregables}, y los apartados siguientes precisan el método, las actividades y las condiciones que permitirán comprobar el cierre de cada fase.
\subsection{Fase metodológica 1: Inicio}\label{sec:inicio}
\subsubsection{Método}\label{sec:inicio-metodo}
La fase de Inicio adaptará DAD para reducir las incertidumbres que condicionan la viabilidad del producto antes de ampliar la producción de escenarios. Se combinarán la revisión de necesidades y antecedentes, la especificación de requisitos, el diseño inicial de arquitectura y las prácticas controladas con FDS, de manera que las decisiones se apoyen en entradas y resultados identificables. Esta adaptación responde a la dependencia del software respecto de datos científicos y recursos de cálculo, por lo que un resultado de fase no será solamente un documento de diseño, sino también evidencia suficiente para preparar y reproducir los casos posteriores. La preparación de esta propuesta constituye un antecedente académico y no sustituye la ejecución de las actividades técnicas que se describen a continuación.

\subsubsection{Actividades [Pendientes: 1. seleccionar sector, 2. confirmar datos y acceso HPC, 3. verificar revisión ejecutable de FDS]}\label{sec:inicio-actividades}
Se consolidarán las necesidades de las entrevistas y se delimitarán requisitos y exclusiones con el equipo y el director, relacionándolos con los objetivos de la sección~\ref{sec:objetivos}. A continuación, se seleccionará el sector a partir de su pertinencia y de la disponibilidad defendible de terreno, coberturas y meteorología, y se preparará un inventario que distinga fuentes identificadas, archivos disponibles, datos transformados y entradas efectivamente comprobadas. Para cada parámetro se conservarán procedencia, unidad y condición de aplicación, evitando completar las brechas con valores de ejemplo. Al mismo tiempo, se describirán mediante los modelos necesarios la organización de componentes, los contratos de intercambio y los metadatos que conectarán la preparación científica con la consulta, dentro de los entregables E1, E2 y E3.

Se realizarán prácticas de FDS que permitan verificar la revisión utilizada, su construcción o instalación y la ejecución de un caso controlado, antes de asumir que el repositorio identificado proporciona una configuración lista para el territorio. También se comprobarán las condiciones de acceso al recurso HPC previsto, ZINE, y las modalidades de ejecución que realmente permita la combinación de simulador e infraestructura. La oferta de recursos y capacitación conocida por el equipo será tratada como antecedente de viabilidad y no como una asignación ya concedida. Finalmente, se prepararán los primeros criterios y casos de prueba del flujo, incluyendo la identificación de condiciones no representadas y la revisión de coherencia espacial, para que las comprobaciones acompañen el desarrollo desde esta fase.

\subsubsection{Resultados esperados [Pendientes: 1. cerrar restricciones de datos y ejecución, 2. aprobar criterios de prueba]}\label{sec:inicio-resultados}
Inicio deberá producir una versión revisada de E1 con requisitos y trazabilidad, una arquitectura inicial E2 y un paquete E3 que documente los datos necesarios recolectados y preparados para el sector acordado, junto con las brechas que limiten su uso. La evidencia de cierre incluirá la reproducción del caso de práctica, la identificación de entradas y salidas, la revisión concreta de WFDS y las condiciones verificadas de acceso y ejecución, además de los criterios iniciales de E6. El equipo y el director deberán poder comprobar que el alcance corresponde a la información disponible y que existe una ruta técnica ejecutable hacia la construcción, por lo que una lista de fuentes sin preparación ni comprobación no bastará para acreditar ese resultado.

La contribución a OE1 será la definición del contrato de consulta y de las condiciones para seleccionar resultados, mientras que OE2 se apoyará en el catálogo previsto y en la viabilidad de ejecución. OE3 recibirá las configuraciones y entradas de práctica reproducibles, y OE4 contará con la arquitectura y las tareas de uso que orientarán su implementación. Estas relaciones no significan que los objetivos queden terminados en Inicio, pues sus productos se completarán mediante los incrementos posteriores. Cualquier restricción que impida sostener un resultado previsto deberá quedar documentada antes de considerar cerrada la fase.


\subsection{Fase metodológica 2: Construcción}\label{sec:construccion}
\subsubsection{Método}\label{sec:construccion-metodo}
Construcción se desarrollará mediante incrementos revisables que combinen trabajo de software y comprobaciones técnicas, con el propósito de integrar gradualmente el flujo de la sección~\ref{sec:solucion}. Cada iteración tendrá un resultado identificable, que podrá consistir en una función de la aplicación o en evidencia necesaria para decidir una configuración, un límite de datos o una condición de consulta. Sobre esa base, se implementará, comprobará y revisará el incremento con el equipo y el director, conservando las versiones evaluadas. Esta adaptación evita suponer que toda investigación sobre simulación debe generar inmediatamente una función de usuario y permite comprobar componentes e integraciones durante la construcción, en lugar de reservar la verificación para la entrega final.

\subsubsection{Actividades [Pendientes: 1. definir casos y tolerancias de contraste, 2. concretar distribución de igniciones y escenarios, 3. fijar reglas de consulta y agregación]}\label{sec:construccion-actividades}
A partir de E3 se preparará un piloto territorial que permita examinar coherencia de entradas, comportamiento numérico y costo de ejecución, registrando tiempo, memoria, configuración y resultados. Se estudiará la sensibilidad de las cantidades de interés frente a los parámetros pertinentes y a la discretización, utilizando rangos defendibles y distinguiendo incertidumbre de entrada de efectos numéricos. Cuando existan observaciones compatibles, el contraste considerará posición del frente en tiempos comparables, dirección de propagación, área o tiempo de llegada, según lo que efectivamente permitan los datos. Un polígono final aislado no se utilizará para demostrar una trayectoria temporal completa. Con esa evidencia se concretarán las configuraciones, la transformación de clases de amenaza en espaciamientos y las combinaciones de condiciones que compondrán E4, sin extrapolar automáticamente el piloto al conjunto de los Cerros Orientales.

La aplicación se construirá integrando almacenamiento y catálogo, obtención meteorológica, búsqueda, tratamiento espacial y animación, conforme a los contratos de E2. Se comprobarán la correspondencia entre consulta y escenarios recuperados, la referencia geográfica, el orden temporal de las salidas y la exclusión del interior quemado. El procedimiento de agregación se contrastará en los casos de referencia que se acuerden para establecer qué diferencias introduce respecto de una representación defendible del incendio, documentando también situaciones desfavorables. Se registrarán por separado el desempeño de la consulta y la calidad de la aproximación, ya que una respuesta rápida no demuestra exactitud física. Las incidencias encontradas se relacionarán con la versión corregida y su nueva comprobación dentro de E6, mientras que el manual E7 se desarrollará a partir del funcionamiento real de los incrementos.

\subsubsection{Resultados esperados [Pendientes: 1. acordar umbrales de desempeño y contraste, 2. fijar cobertura mínima evaluable]}\label{sec:construccion-resultados}
El resultado de Construcción será una versión integrada E5 lista para evaluación con usuarios, acompañada de la biblioteca E4, el paquete reproducible E3 actualizado, la arquitectura E2 y los registros de E6 correspondientes a los incrementos. La condición de cierre será poder ejecutar el flujo completo bajo condiciones identificadas y comprobar, mediante pruebas vinculadas a requisitos, que cada componente intercambia e interpreta correctamente sus datos. Además, deberá existir un informe del piloto que permita reconocer las condiciones examinadas, las diferencias de aproximación y las restricciones restantes. Los criterios cuantitativos de cobertura, latencia y contraste se acordarán antes de evaluar su cumplimiento y no se reemplazarán por una apreciación visual de las animaciones.

OE1 quedará respaldado por consultas reproducibles y por la identificación de sus resultados, OE2 por los casos ejecutados y catalogados dentro de la cobertura acordada, y OE3 por configuraciones defendibles con evidencia de preparación y contraste. OE4 contará con la integración georreferenciada y las primeras comprobaciones de interacción. El cierre técnico deberá distinguir lo que funciona como prototipo exploratorio de lo que dispone de validación territorial, de manera que una demostración del software no se presente como habilitación para uso operativo. Si la agregación solamente resulta defendible bajo ciertas condiciones, esas condiciones deberán acompañar el producto evaluable.


\subsection{Fase metodológica 3: Transición}\label{sec:transicion}
\subsubsection{Método [Pendientes: 1. definir protocolo y criterios de evaluación con usuarios]}\label{sec:transicion-metodo}
Transición adaptará DAD a la evaluación de una versión demostrable y a su entrega académica, utilizando tareas de uso que permitan observar cómo el usuario interpreta los resultados y sus restricciones. La evaluación con bomberos o personas que representen el rol previsto se organizará sobre casos y condiciones explícitas, sin confundir participación en una sesión con aceptación institucional o autorización operativa. Los resultados de las pruebas técnicas de Construcción servirán como base y se complementarán con evidencia de comprensión, de modo que la valoración del producto no dependa únicamente de que la interfaz sea navegable o la animación se reproduzca correctamente.

\subsubsection{Actividades [Pendientes: 1. confirmar participantes y tareas de evaluación]}\label{sec:transicion-actividades}
Se prepararán tareas representativas del flujo existente y se definirán las condiciones de aplicación, los resultados esperados y la forma de registrar observaciones. Durante las sesiones se examinará si los participantes logran plantear una consulta, reconocer la información espacial y temporal presentada e identificar los límites de la estimación, evitando solicitar acciones que excedan las funciones comprometidas. Los hallazgos se incorporarán al informe E6 y se tratarán mediante correcciones verificables de la versión evaluada. Después de cada ajuste pertinente se repetirán las comprobaciones afectadas, conservando la relación entre problema observado, modificación y resultado de la nueva prueba.

Se consolidarán el código y el ejecutable E5, las configuraciones E3 y el catálogo E4 necesarios para reproducir la demostración, y se revisará la correspondencia final con E1 y E2. El manual E7 describirá instalación, dependencias, uso del flujo, condiciones representadas y límites conocidos, mientras que el informe final reunirá resultados y conclusiones sustentados por las evidencias obtenidas. La entrega al director incluirá una demostración preparada con datos y versiones identificadas, de manera que pueda distinguirse lo ejecutado de las funciones o condiciones que hayan quedado fuera de la evaluación.

\subsubsection{Resultados esperados [Pendientes: 1. acordar criterios finales de aceptación académica]}\label{sec:transicion-resultados}
La fase deberá cerrar con una versión demostrable corregida, el informe de evaluación E6, el manual y el informe final E7, junto con la matriz de trazabilidad E1 y los artefactos técnicos actualizados que respaldan la entrega. Su cumplimiento se comprobará mediante la reproducción de la demostración, la revisión de resultados frente a los criterios acordados y la identificación explícita de limitaciones o problemas no resueltos. La documentación de una evaluación desfavorable también constituirá evidencia del proyecto y deberá reflejarse en sus conclusiones, en lugar de omitirse para presentar un resultado aparentemente satisfactorio.

La evaluación de comprensión completará principalmente la evidencia de OE4, mientras que la demostración integrada y la trazabilidad final permitirán revisar conjuntamente los aportes de OE1, OE2 y OE3. Esta fase no sustituye las pruebas realizadas durante el desarrollo, sino que consolida su evidencia en una entrega evaluable. La aceptación académica se determinará con el director a partir de los criterios y evidencias consolidados.



\clearpage
\section{Aspectos generales del proyecto}\label{sec:aspectos}
\subsection{Compromiso de apoyo de la institución}\label{sec:apoyo}
No aplica. El proyecto se dirige a la comunidad y a quienes encuentren útil la aplicación, especialmente personas vinculadas a la atención de incendios, y no se desarrolla por encargo de una institución cliente ni se ha identificado una entidad receptora con un compromiso formal de apoyo. En consecuencia, no se incorpora una carta institucional a esta propuesta. La disponibilidad prevista de infraestructura universitaria para apoyar la ejecución técnica se aborda en la sección~\ref{sec:inicio-actividades} y no constituye, por sí misma, un compromiso de una institución usuaria del producto.

\subsection{Derechos patrimoniales [Pendientes: 1. definir la declaración correspondiente]}\label{sec:derechos}


\clearpage
\section{Marco teórico}\label{sec:marco}
\subsection{Fundamentos y conceptos relevantes para el proyecto}\label{sec:fundamentos}
\textbf{Modelado físico del incendio.} FDS es el software de simulación que representa el incendio mediante la interacción entre combustible, transferencia de calor, combustión y movimiento del aire. El trabajo de Mell y colaboradores constituye un antecedente de la aproximación física a incendios de pastizales~\cite{mell2007}, mientras que Valero, Jofre y Torres estudian el uso de configuraciones de distinta fidelidad para analizar propagación, sensibilidad e incertidumbre~\cite{valero2021}. En su formulación CFD, FDS utiliza una representación de las ecuaciones de movimiento para flujos de bajo número de Mach impulsados por calor, por lo cual resuelve un problema tridimensional y transitorio. Las capacidades de representación de vegetación asociadas históricamente a WFDS se utilizan dentro de FDS según la convención de la sección~\ref{sec:solucion}~\cite{fdsnist,wfdsusfs}. Esta formulación se distingue de los modelos de propagación de frente.

\textbf{Discretización y costo de cálculo.} La malla CFD divide el dominio físico en celdas donde se aproximan las variables del flujo y sus intercambios. Si un dominio rectangular se discretiza de manera uniforme, su número de celdas puede expresarse como $N=(L_x/\delta x)(L_y/\delta y)(L_z/\delta z)$, donde las longitudes y los tamaños de celda deben utilizar unidades compatibles. Por tanto, reducir a la mitad el tamaño de celda en las tres direcciones multiplica geométricamente por ocho el número de celdas del mismo dominio, aunque el costo temporal real también depende de la integración, los modelos y la ejecución. Este razonamiento explica por qué la fidelidad requiere decisiones explícitas de resolución, pero no proporciona una medición de tiempo o memoria para el sector del proyecto. El análisis de Valero y colaboradores aporta un antecedente para estudiar esa relación y sus efectos en cantidades de interés~\cite{valero2021}.

\textbf{Combustibles y representación territorial.} Una cobertura vegetal cartográfica identifica clases y distribución espacial, pero no define por sí sola la estructura tridimensional ni todas las propiedades físicas que requiere una entrada de simulación. La preparación debe distinguir composición y geometría, carga, humedad y las propiedades pertinentes al modelo elegido, conservando sus unidades y condiciones de aplicación. En el contexto de los Cerros Orientales, Verano Velásquez caracteriza información territorial para FARSITE~\cite{verano2013}, y Méndez Álvarez estudia la calibración de modelos de combustible altoandino~\cite{mendez2018}. Sus aportes sirven para examinar fuentes, parámetros y dificultades locales, aunque la transferencia a FDS exige comprobar correspondencias y restricciones. Una calibración realizada para otro modelo no se convierte automáticamente en una calibración de FDS, ni una cobertura histórica equivale a un inventario actual y exacto de especies y combustibles.

\textbf{Amenaza y distribución espacial de igniciones.} Morales Salguero integra información territorial mediante álgebra de mapas y simulación para estudiar amenaza de ignición y propagación en los Cerros Orientales~\cite{morales2017}. Las categorías resultantes permiten distinguir niveles de amenaza, pero no deben interpretarse como probabilidades estadísticas calibradas de que ocurra una ignición. Esta distinción es relevante para fundamentar la distribución de referencia establecida en la sección~\ref{sec:solucion}, cuya regla de espaciamiento debe ser definida por el proyecto. Los puntos de ignición representan posiciones donde se preparan escenarios, mientras que las celdas CFD representan la discretización de cada experimento físico, de modo que aumentar puntos disponibles no equivale a refinar la resolución de las ecuaciones.

\textbf{Cómputo de alto desempeño y experimentos precalculados.} Existen dos niveles de paralelización que requieren comprobaciones diferentes: ejecutar escenarios independientes de forma concurrente y distribuir el cálculo de un escenario entre recursos que intercambian información. El primer nivel aprovecha la independencia de combinaciones de entrada, mientras que el segundo depende de la descomposición y comunicación admitidas por la revisión del simulador. La distribución de escenarios independientes y la distribución interna de un caso plantean, por tanto, necesidades diferentes de memoria, almacenamiento y comunicación. Esta distinción fundamenta las comprobaciones de viabilidad y las mediciones del piloto descritas en las secciones~\ref{sec:inicio-actividades} y~\ref{sec:construccion-actividades}.

\textbf{Consulta por semejanza.} Recuperar escenarios requiere representar la situación observada y los casos disponibles mediante variables comparables, con una definición explícita de cercanía y tratamiento de unidades. Cover y Hart ofrecen un antecedente fundacional sobre vecinos cercanos en clasificación~\cite{cover1967}, del cual puede retomarse el principio de utilizar proximidad entre representaciones para recuperar información. Sin embargo, sus resultados de clasificación no demuestran la precisión física de una búsqueda de incendios ni seleccionan automáticamente la métrica adecuada para este problema. La normalización de variables, la dirección del viento, la proximidad espacial y las condiciones no cubiertas deberán tratarse según las propiedades de la biblioteca, y la cercanía de un resultado deberá comunicarse como correspondencia entre condiciones, sin convertirla en garantía de representatividad.

\textbf{Incertidumbre, multifidelidad y agregación.} La incertidumbre puede proceder de datos de entrada, simplificaciones del modelo, discretización y diferencias entre las condiciones observadas y las disponibles. Monte Carlo estudia variabilidad mediante repeticiones con entradas muestreadas, mientras que los enfoques multinivel combinan estimaciones de distinta fidelidad para reducir el costo de estudiar ciertas cantidades estadísticas. El artículo seminal de Valero y colaboradores utiliza ese marco para analizar simulación de incendios~\cite{valero2021}. Su objeto es diferente del procedimiento espacial de la sección~\ref{sec:solucion}: la estimación estadística entre fidelidades no constituye una regla de combinación de frentes. La superposición de frentes procedentes de igniciones calculadas de forma independiente no reproduce necesariamente todas las interacciones de un incendio iniciado desde un perímetro, por lo que requiere un contraste propio. En consecuencia, multifidelidad y Monte Carlo multinivel se conservan como fundamentos bibliográficos y no se incorporan por esta referencia como componentes comprometidos del producto.

\textbf{Representación geográfica e interpretación.} Una visualización georreferenciada vincula los resultados a un sistema espacial y a una secuencia temporal coherentes con la información consultada. Su corrección exige comprobar coordenadas, unidades, orientación y correspondencia de tiempos, además de diferenciar lo observado de lo estimado. Estas propiedades permiten examinar que la animación represente los datos procesados, aunque no bastan para acreditar que el usuario comprenda su alcance. Por ello, la evaluación distinguirá la coherencia técnica de la representación y la comprensión de las tareas de uso, conforme a las fases de las secciones~\ref{sec:construccion} y~\ref{sec:transicion}.

\subsection{Alternativas de solución [Pendientes: 1. establecer tres alternativas adicionales, 2. documentar sus alcances e impactos, 3. elaborar la tabla comparativa con la propuesta]}\label{sec:alternativas}


\clearpage
\refstepcounter{section}\label{sec:referencias}
\addcontentsline{toc}{section}{6 Referencias}
\renewcommand{\refname}{6 Referencias}

\bibliographystyle{ieeetr}
\begingroup
\small
\setlength{\parskip}{0pt}
\bibliography{references}
\endgroup
\end{document}
